\chapter{Starea actual\u a a domeniului}
\vspace*{-1cm}
\^In acest capitol este analizat contextul \^in care se \^incadreaz\u a aplica\c tia web pentru organizarea \c si gestionarea evenimentelor: cercet\u ari relevante pentru comer\c tul electronic, ticketing online \c si securitate web, platforme existente \c si o compara\c tie func\c tional\u a.

\section{Studii \c si articole \c stiin\c tifice}

\subsection{Arhitecturi web \c si API-uri REST}

Fielding \cite{fielding2000} define\c ste principiile REST (\textit{Representational State Transfer}) ca stil arhitectural pentru sisteme distribuite. Separarea clientului de server, utilizarea resurselor adresabile prin URI \c si a metodelor HTTP standard (GET, POST, PUT, DELETE) faciliteaz\u a dezvoltarea aplica\c tiilor web scalabile. ManFast expune un API REST \^in Node.js/Express, consumat de frontend-ul React abordare aliniat\u a cu recomand\u arile din literatura de specialitate.

Fowler \cite{fowler2002} descrie tipare de proiectare enterprise, inclusiv separarea responsabilit\u a\c tilor \^intre straturi (prezentare, logic\u a de business, persisten\c t\u a). Aplica\c tia analizat\u a respect\u a acest principiu: React gestioneaz\u a UI, Express con\c tine regulile de business, PostgreSQL stocheaz\u a datele.

\subsection{Securitatea autentific\u arii}

OWASP \cite{owasp2024} enumer\u a bune practici pentru autentificare: hash-uirea parolelor (nu stocarea \^in clar), protec\c tia sesiunilor, limitarea ratei de cereri, validarea input-ului. ManFast implementeaz\u a bcrypt pentru parole, JWT cu expirare 24h, rate limiting pe codurile de verificare a e-mailurilor \c si resetare parolei.

Standardul JWT (JSON Web Token) \cite{jwt2015} permite transmiterea securizat\u a a claims-urilor \^intre p\u ar\c ti. Token-ul semnat con\c tine id-ul utilizatorului \c si rolul (user/admin), verificat la fiecare cerere protejat\u a.

Provos \c si Mazi\`eres \cite{bcrypt1999} introduc bcrypt, algoritm de hash adaptat la cost computa\c tional configurabil, rezistent la atacuri brute-force pe parole. ManFast folose\c ste cost factor 10.

\subsection{Pl\u a\c ti online \c si gateway-uri}

Documenta\c tia Stripe \cite{stripe2024} descrie fluxul Checkout: sesiune de plat\u a g\u azduit\u a, redirect c\u atre pagina securizat\u a Stripe, confirmare prin webhook sau verificare sesiune. ManFast creeaz\u a sesiuni Checkout cu metadata (event\_id, user\_id, nume participan\c ti) \c si confirm\u a plata prin endpoint dedicat, f\u ar\u a a stoca date de card pe serverul propriu --- reduc\^and responsabilitatea PCI-DSS.

\subsection{Persisten\c t\u a \c si baze de date rela\c tionale}

PostgreSQL \cite{postgresql2024} ofer\u a tranzac\c tii ACID, constr\u angeri de integritate referen\c tial\u a \c si tipuri \linebreak avansate (ex.: array pentru imagini eveniment). Schema ManFast exploateaz\u a aceste capabilit\u a\c ti pentru consisten\c ta bilete/locuri disponibile.

\subsection{Interfa\c t\u a utilizator \c si SPA}

React \cite{react2024} \c si ecosistemul s\u au permit construirea de \textit{Single Page Applications reactive}. Material UI \cite{mui2024} ofer\u a componente accesibile \c si responsive. ManFast utilizeaz\u a React 19, Vite pentru build rapid \c si MUI 9 cu tem\u a dark personalizat\u a.

\section{Aplica\c tii dedicate domeniului}

\subsection{Eventbrite}

Eventbrite \cite{eventbrite2024} este o platform\u a interna\c tional\u a pentru crearea \c si promovarea evenimentelor, cu v\^anzare bilete, pagini de destina\c tie \c si instrumente de marketing. Ofer\u a check-in prin cod QR \c si analiz\u a audien\c t\u a. Platforma suport\u a multiple tipuri de evenimente: concerte, conferin\c te, workshop-uri. Fluxul tipic: organizatorul creeaz\u a pagina evenimentului, seteaz\u a pre\c turile \c si capacitatea, promoveaz\u a link-ul, iar participan\c tii achizi\c tioneaz\u a bilete online.

Limit\u ari pentru organizatori mici: comisioane per bilet, personalizare redus\u a a fluxului de check-in, dependen\c t\u a de ecosistemul proprietar. Datele participan\c tilor sunt stocate pe infrastructura Eventbrite, nu \^in controlul direct al organizatorului. Pentru conferin\c te cu sesiuni paralele sau concerte cu locuri numerotate, unele func\c tionalit\u a\c ti avansate pot lipsi sau necesit\u a addon-uri pl\u atite.

Din perspectiva tehnologic\u a, astfel de platforme folosesc arhitecturi cloud scalabile, procesare pl\u a\c ti prin PSP integrat \c si aplica\c tii mobile pentru scanare QR. ManFast adopt\u a un subset al acestor capabilit\u a\c ti, orientat spre simplitate \c si control local al datelor \^in PostgreSQL.




\subsection{Platforme locale \c si generice}

Pe pia\c ta rom\^aneasc\u a exist\u a solu\c tii pentru bilete la concerte \c si festivaluri (iaBilet, Eventim), precum \c si platforme e-commerce generice (WooCommerce, Shopify) care pot vinde ``produse'' reprezent\^and locuri la eveniment. Acestea sunt optimizate pentru v\^anzare la scar\u a larg\u a, nu neap\u arat pentru meeting-uri corporate sau conferin\c te universitare.

Organizatorii locali folosesc frecvent formulare Google, pl\u a\c ti directe \c si liste Excel solu\c tii cu cost zero ini\c tial, dar cu risc ridicat de erori \c si lips\u a de audit trail pentru pl\u a\c ti. Reconcilierea manual\u a a listei de participan\c ti cu transferurile bancare consum\u a timp \c si poate genera dispute la intrarea \^in sal\u a.

Aplica\c tia propus\u a \^incerc\u a s\u a ofere un echilibru: cost de operare redus (self-hosted), flux clar de la descoperire eveniment la plat\u a \c si check-in, f\u ar\u a complexitatea unei platforme glo-bale. Adaptarea la jude\c tele rom\^ane\c ti \c si pl\u a\c ti RON via Stripe reprezint\u a un avantaj competitiv \^in contextul local fa\c t\u a de solu\c tiile generice interna\c tionale.

\subsection{Tendin\c te \c si standarde \^in industria ticketing-ului}

Industria biletelor electronice a crescut accelerat dup\u a pandemia COVID-19, c\^and contactless check-in \c si pl\u a\c tile online au devenit norm\u a. Platformele mature investesc \^in trei direc\c tii: (1) experien\c t\u a mobil\u a — aplica\c tii native sau PWA pentru scanare rapid\u a; (2) analitic\u a — dashboard-uri de v\^anz\u ari, heatmap-uri de pre\c turi dinamice; (3) integrare marketing — email automation, remarketing pe re\c tele sociale.

ManFast nu urm\u are\c ste replicarea tuturor acestor capabilit\u a\c ti enterprise, ci func\c tionalitatea minim\u a viabil\u a pentru organizatori locali: publicare eveniment, v\^anzare online, list\u a \linebreak participan\c ti, check-in la intrare. Aceast\u a decizie de scope este aliniat\u a cu principiul MVP (\textit{Minimum Viable Product}) din literatura de product management.

\subsection{Conformitate PCI-DSS \c si delegarea pl\u a\c tilor}

Standardul PCI-DSS (\textit{Payment Card Industry Data Security Standard}) impune m\u asuri stricte pentru orice sistem care proceseaz\u a, stocheaz\u a sau transmite date de card. Solu\c tiile care folosesc Stripe Checkout hosted deleg\u a colectarea datelor sensibile c\u atre Stripe comerciantul prime\c ste doar confirmarea pl\u a\c tii prin API. \newpage ManFast adopt\u a explicit acest model: frontend-ul redirec\c tioneaz\u a utilizatorul pe domeniul \textbf{checkout.stripe.com}, iar backend-ul verific\u a statusul sesiunii f\u ar\u a a vedea num\u arul de card.

Aceast\u a abordare reduce suprafa\c ta de atac \c si simplific\u a auditul de securitate pentru o lucrare de licen\c t\u a, men\c tin\^and totodat\u a un flux profesional de plat\u a recunoscut de utilizatori.

\subsection{Arhitectura SPA \c si API REST}

Aplica\c tiile \textit{Single Page Application} (SPA) [8] \^incarc\u a o singur\u a pagin\u a HTML \c si navigheaz\u a prin JavaScript, ob\c tin\^and date de la server prin API JSON.

 \textbf{Avantaje}: interfa\c t\u a reactiv\u a, separare clar\u a frontend/backend, posibilitatea de a extinde ulterior cu aplica\c tie mobil\u a consum\^and acela\c si API.

\textbf{Dezavantaje} relevante pentru ManFast: SEO limitat pe pagini dinamice (acceptabil pentru platform\u a cu autentificare), necesitatea gestion\u arii token-ului JWT pe client. Compara\c tia cu arhitectura server-side rendering (Next.js) ar fi o extensie viitoare; pentru licen\c t\u a, Vite + React ofer\u a timp de dezvoltare redus \c si ecosistem bogat (Material UI).

\subsection{Studiu de caz: fluxul utilizatorului \^in platforme consacrate}

Pe Eventbrite, utilizatorul parcurge: descoperire eveniment → selectare tip bilet → checkout → email confirmare cu PDF/QR. Organizatorul acceseaz\u a dashboard cu statistici v\^anz\u ari \c si export CSV participan\c ti.

ManFast simplific\u a acest flux pentru organizatori locali (conferin\c te, concerte, meeting-uri):

- un singur tip de bilet pe eveniment (pre\c t unic);

- nume nominal obligatoriu (identificare la check-in);

- f\u ar\u a marketplace global, evenimentele sunt listate doar pe instan\c ta ManFast;

- check-in manual din panou admin, f\u ar\u a aplica\c tie mobil\u a dedicat\u a scan\u arii QR (extensie viitoare).

Aceast\u a simplificare reduce complexitatea implement\u arii \c si timpul de \^inv\u a\c tare pentru administratori neexperimenta\c ti \^in IT.

\section{Criterii de evaluare a solu\c tiilor}

Pentru a fundamenta compara\c tia din \textit{Tabelul 2.2}, au fost definite criterii ponderate:

Adecvare func\c tional\u a (40\%) — suport bilete nominale, check-in, filtrare local\u a.

Cost total de operare (20\%) — comisioane PSP, abonamente platform\u a, hosting.

Control date (20\%) — unde sunt stocate datele participan\c tilor, export, GDPR.

Personalizare (10\%) — posibilitatea modific\u arii fluxului sau branding.

Securitate autentificare (10\%) — verificare email, reset parol\u a, roluri.

ManFast ob\c tine scor ridicat la control date \c si personalizare (cod deschis, PostgreSQL propriu), scor mediu la cost (Stripe per tranzac\c tie, hosting self-managed), \c si scor bun la adecvare func\c tional\u a pentru ni\c sa aleas\u a.

Aceste analize suplimentare contextualizeaz\u a pozi\c tionarea ManFast fa\c t\u a de solu\c tiile \newline  existente \c si justific\u a alegerile de scope ale proiectului.

\section{Analiz\u a comparativ\u a}

Tabelul~\ref{tab:comparatie} sintetizeaz\u a diferen\c tele \^intre platforme reprezentative \c si aplica\c tia propus\u a.

\begin{table}[htbp]
	\centering
	\renewcommand{\arraystretch}{1.5} 
	\setlength{\tabcolsep}{6pt} 
	\begin{tabular}{|>{\raggedright\arraybackslash}p{3.5cm}|>{\raggedright\arraybackslash}p{2.2cm}|>{\raggedright\arraybackslash}p{2.5cm}|>{\raggedright\arraybackslash}p{3.5cm}|}
		\hline
		\textbf{Func\c tionalitate} & \textbf{Eventbrite} & \textbf{Generic e-comm} & \textbf{Aplica\c tia propus\u a} \\ \hline
		Bilete nominale & Da & Par\c tial & Da \\ \hline
		Pl\u a\c ti Stripe RON & Nu (alte PSP) & Variabil & Da \\ \hline
		Check-in admin & Da (QR) & Nu & Da (list\u a) \\ \hline
		Verificare email register & Da & Variabil & Da (cod 6 cifre) \\ \hline
		Filtrare jude\c t/dat\u a & Da & Nu & Da (42 jude\c te) \\ \hline
		Cod deschis / personalizabil & Nu & Nu & Da \\ \hline
		Reset parol\u a tel.+email & Variabil & Variabil & Da \\ \hline
	\end{tabular}
	\caption{Compara\c tie func\c tional\u a}
	\label{tab:comparatie}
\end{table}

\section{Concluzii par\c tiale}

Analiza demonstreaz\u a c\u a exist\u a un spa\c tiu pentru o solu\c tie adaptat\u a: conferin\c te, meeting-uri, concerte \c si evenimente culturale, cu pl\u a\c ti RON via Stripe, bilete pe nume \c si admi-nistrare simpl\u a. Aplica\c tia propune o implementare open-source, extensibil\u a, care integreaz\u a func\c tionalit\u a\c tile esen\c tiale \c si le aliniaz\u a cu cerin\c tele locale (jude\c te, SMTP Brevo pentru e-mailuri).

\vspace{0.3cm}
\noindent Urm\u atorul capitol detaliaz\u a arhitectura, tehnologiile \c si implementarea solu\c tiei propuse.